%\chapter{Future Directions for Agentic LLMs}\label{chap:concl}
\chapter{Future Directions for Agentic LLMs}\label{chap:concl}


Having arrived at the end of this book, we provide a brief high-level summary of the material that we have discussed. First we look back at our chapters about how agentic LLMs are built. Then we summarize our discussion of different applications of agentic LLMs and their behavior in single and multi-agent settings. Finally, in Section~\ref{sec:limitations} we look at the limitations of current approaches, and what the future may bring.

\section{The State of Agentic LLMs}
As discussed in Parts I, II, and III, the foundations of agentic LLMs lie in language modeling, reasoning, tool use, and different degrees of autonomy.

\subsection{Technology} 
Agentic LLMs are reasoning  models that act in the world and adapt their behavior from feedback. After having been pretrained and finetuned (Chapter~\ref{chap:pipeline}), at inference time they may access to external information (retrieval augmented generative AI, Chapter~\ref{sec:rag}). They may also prompt themselves at inference time to verify results, leading to self-reflection and in-context reinforcement learning, Chapter~\ref{chap:reasoning}. Agentic LLMs have become multi-modal, they operate in language, vision, and action domains (Chapter~\ref{chap:world}). To do so effectively, agentic LLMs can be provided with pretrained world models with a rudimentary understanding of the physical world. Through interacting with the environment, these world models may be improved at test time. 

The main machine learning architecture used is the transformer architecture. Models are pretrained through self-supervision (next-token prediction) and finetuned on specific tasks through supervised and reinforcement learning. 

In-context chain of thought prompting has improved reasoning abilities of models significantly. Inference-time methods have been developed for finetuning (such as RLVR and GRPO) and self-reflection (such as ReAct and reflexion, Chapter~\ref{chap:reasoning}). A whole spectrum of agent-environment interaction prompting styles have been developed (Chapter~\ref{chap:agents}).

\subsection{Applications}
In Part IV, various tools have been introduced to ease the construction of agentic LLMs (Chapter~\ref{chap:tools}). Agentic LLMs have been designed for many different applications. In e-commerce, they can help understand the preferences of their users and act on recommendations. In  medicine, they can read and summarize diagnostic texts. In finance, they can analyze financial reports and provide recommendations to traders. In science and education, they can help in literature research, they can generate scripts to run experiments, they can analyze results, they can help write papers, and they can help review them (Chapter~\ref{sec:assist}). 
%(Niki: would be good to add references explicit to these parts.

\subsection{Behavior}
In Part V we shifted gears from looking at agentic LLMs through a technical lens, to studying how they behave in single-agent and multi-agent settings. After discussing basic behavioral features of LLMs, in Chapter \ref{chap:behavior} we discussed how LLM-driven agents perform on tasks adapted from human behavioral approaches and frameworks developed especially for AI benchmarking. Chapter \ref{chap:role} looked closely at multi-agent interactions that built on predefined roles and protocols. In Chapter \ref{chap:emergent} focus shifted to open-ended settings in which the object of study were the behavioral structures, norms, and conventions that emerged at system level from small-scale interactions between agents.

Various recurring patterns were identified throughout the approaches and studies discussed in Part V. These include the discussions on construct validity (how do observations generalize to the human situation), prompt sensitivity and cascading effects (especially in multi-agent settings), the knowing-doing gap (models' capabilities at verbal predicting often did not translate well into actions), positive effects of structured reasoning (increasing reasoning budgets was particularly useful when structured along lines inspired on human reasoning), and emergent properties in desired directions (performance above baseline matched for compute budget) and undesired directions (norm violations) were observed.

\section{Limitations and Future Work}\label{sec:limitations}
%\section{Conclusion}
Agentic LLMs build on the power of large language models, and have greatly expanded the capabilities and uses of LLMs. However, there are still many limitations, which we will discuss next.
%Agentic LLMs have shown how the capabilities of language models can be expanded by interacting with the world around them. We have discussed their  impressive achievements. However, limitations remain that offer opportunities for further research.
We progress from memory, scale, explainability, and cognitive architectures, which are more technical, to topics with broad implications such as potential sentience or consciousness and ethical implications that follow from such properties. For some of these topics, solution directions have been discussed in the preceding chapters. 

\subsection{Memory and Self-reflection}
Language models are functions: the output that they compute is only  based on their input---they are memory-less between calls. In order for agents to act naturally in the world, to hold a conversation, to be able to learn from their actions by self-reflection, agents must be able to remember. For us, our identity hinges on our ability to remember. Agents require a form of memory.

A simple way to achieve a rudimentary form of memory is  to feed the past interaction history back into the prompt---and current chatbots use this method with success. However, context windows remain limited in size, despite some success with research  into long-context models. Agentic memory methods are an active area of research. LLMs should learn how to best compact information and store it in long or short term memory.

Agents can access outside information through retrieval-augmented generative AI methods (RAG), for example from corporate databases. 
Elaborate agentic memory architectures are being developed to provide agents with longer term memory, either textual, or token embeddings, in vector databases  (Section~\ref{sec:memory}). 
%Reflexion has short term and long term memory.

\subsection{Scale}
As agents become able to self-reflect and to learn from interactions with their environment, we should discuss the topic of scaling of the underlying language models.

A major challenge in the development of large language models is how to further improve their performance. Originally, size was the major determinant---hence the name {\em large} language models---only models of sufficient size and trained on sufficiently large text corpuses achieved adequate benchmark performance, and only large models exhibited in-context learning.  Scaling the size of the model and of the data thus appeared to be the way forward. However, scaling further is becoming more and more difficult, as most high-quality data that is available on the Internet has been now used to train current models. Also, training longer and larger models is expensive; time, money and power consumption are limiting this direction. Current methods are reaching a plateau \citep{sutskever2024}.

To improve LLM performance further, it is important to continue to acquire  training
data that is  challenging (Chapter~\ref{chap:synthetic-data}). One approach is to use the interactions of agentic LLMs with their environment as training data. When this interaction data is new, diverse, and high quality, it can be used for pretraining (see also Figure~\ref{fig:virtuous}). Note that already inference-time methods such as GRPO and RLVR are using inference-time results for finetuning with good results \citep{guo2025deepseek}, and VLAs are  trained on self-generated action data \citep{black2024pi_0,brohan2023rt}, Chapter~\ref{chap:world}. 

\paragraph{Overthinking}
Scaling studies into chain of thought also reveal limitations of more thinking effort. Chain of thought works better for some domains than for others \citep{sprague2024cot}. Furthermore, deeper reasoning may lead to {\em overthinking}, where the reasoning process becomes less effective as more effort is spent \citep{liu2024mind,wu2025more}. Studies are ongoing and reveal more insight into how chain of thought works \citep{chen2025does}, and which mechanism are effective \citep{chen2026think}.

\paragraph{Open World Models and Curriculum Learning}
For the scaling approach to work it is important that synthetic training data is diverse, and different from the existing data. Open world models are studied to specifically generate diverse training environments. The idea of an open world-model is that it can create infinite datasets or environments, in which agents can continue to learn, to keep improving their intelligence \citep{bruce2024genie}. Open world models are inspired by evolutionary algorithms, quality diversity, and computer games \citep{vinyals2019grandmaster,mouret2015illuminating,meta2022human}.

Another approach to further scale model performance is to create better training methods. Curriculum learning is built on the observation that most learning works better if a system is first trained on easy tasks, before harder tasks are tried. (This holds both for human brains and for artificial neural networks.) In most domains a hierarchy of concepts exists that are related and build upon each other. (In school we learn addition before we learn long division, or even logarithms.) Curriculum methods are being tried for both pre- and posttraining \citep{chen2025self,elgaar2026curriculum}.

\subsection{Hallucination and Explainability}
Another  limitation of current agentic LLMs is the problem of hallucination: models often provide answers that sound plausible but are factually incorrect, and their stated reasoning may not be causally responsible for the output they produce---a property known as \emph{faithfulness}. 

The field of explainable AI (XAI) provides a spectrum of techniques to address this, ranging from behavioral and attribution-based methods such as SHAP and LIME, which treat the model as a black box, to concept-based interpretability, which probes internal representations for human-understandable features. Going further, mechanistic interpretability (MI), as discussed in Chapter~\ref{chap:mi}, aims to reverse-engineer the computational circuits inside transformer models. 

Tools such as sparse autoencoders, neural lenses and activation patching have already been used to identify specific circuits responsible for hallucination and to provide empirical evidence that chain of thought reasoning can be a faithful, causally active component of model behavior in sufficiently large architectures. As agentic LLMs are increasingly deployed in high-stakes settings, MI offers a path toward hallucination detectors that trigger when a model's internal confidence features diverge from its external output---a critical step toward accountable agentic systems.

A step further is the comparison of cognitive architectures between humans and transformers. Language models appear to produce natural language in a way that is similar to humans. Is that because the architectures that produce the utterances are similar, is it by coincidence, or is there another explanation? Chapter~\ref{chap:cognition} explored this question.
%\todo{Max; please check and add}

\subsection{Emotional Behavior and Consciousness}
%- anthropomorphization and consciousness: attributing human elements to non-human entities, including when an entity should be considered conscious.
Humans have a tendency to anthropomorphize: to attribute purposeful behavior to inanimate objects in their surroundings, when they behave in ways that are similar to our behavior, and sometimes even when they do not. Implicitly, we may attribute consciousness to such objects. This leads us to believe and say things such as {\em the computer is thinking}, {\em the computer wants}, and {\em the computer must not be feeling well} (or even {\em this book loves to be read}). LLMs suffer the same fate, and their utterances may from time to time persuade even the most rational philosophers to believe that they have a purpose (whereas they clearly cannot have a purpose, since they are merely algorithms that are being performed by a mechanical device---unlike humans, since we are organic).\footnote{Note that the duck-typing point of view disagrees, this behaviorist view holds that if it walks like a duck and quacks like a duck, then we consider it a duck.} 
% er is nog een variant die niet op het wat infantiele idee van duck-typing berust: intelligence/life is computational zoals Blaise AyA het verdedigt. LLMs zijn computerprogramma's, wij zijn computerprogramma's; de hardware en software verschillen op veel punten, maar niet op alle punten. Ik ga proberen dit in 21 helder te krijgen zonder dat het veel pagina's worden.

Reports have been published about emotional relations between chatbots and humans, sometimes leading to dangerous behavior, and worse. 

The topic of consciousness is actively debated by AI philosophers. Consciousness is generally assumed to mean being aware of oneself and the environment, and having subjective experiences, thoughts, and feelings \citep{velmans2009understanding,butlin2023consciousness}. But how do we test for consciousness? 
A pragmatic view is that for our society, behavior is what counts. If an entity acts, the consequences of its actions are what matters, and less so whether the entity was conscious or not. \citet{bengio2025illusions} argue that the mere illusion of consciousness, through actions, has consequences, and that it will become increasingly difficult for us to discern  conscious entities from non-conscious entities. They argue that society must find answers on how to deal with agents that act, whether in medicine, finance, or science. Both credit and liability must ultimately be attributed in a fair way, and answers must be found on how to handle self-preserving AI systems. %\todo{Michiel?} 

\subsection{Social Behavior}
We have discussed examples of how LLMs can simulate reasoning by humans (Chapter~\ref{chap:behavior}), and how LLMs can improve simulations of social behavior in multi-agent systems. A limitation of current simulations of social behavior of agents is that the behavior is heavily influenced by the  prompt script.  Although some studies found realistic social behavior, many studies also found limits where the LLMs fail.  

Furthermore, the study of emergent behavior at scale poses challenges (Chapter~\ref{chap:emergent}). 
As with all multi-agent simulations, validation of the emergent behavior, to see how realistic  the simulated behavior is, is a challenge, although some studies report impressive results \citep{park2024generative}. %\todo{Max: uitbreiden?}   

\subsection{Legal and Ethical Aspects}
%to act in the real wor- legal: explainable, responsible, liable, ethical behavior by algorithms. Beyond our current expertise, deserves dedicated book.
Actions by assistants in the real world have real world consequences. Many legal and ethical challenges exist for  agentic LLMs that act, when the actions of imperfect models have real world consequences. When a financial trading assistant hallucinates, or when a self driving robot makes a wrong inference, questions of responsibility and liability should be addressed, before such systems are deployed in the real world. 

%An interesting trade-off occurs between breadth of the application domain and accuracy of the answer.
%The application areas  for the assistants in this survey---shopping, medical diagnosis, finance---are  narrow. The narrower the application domain, the better the answers. 

%More legal and ethical questions arise, for example, on privacy and fairness, and, possibly, concerning the rights of algorithmic entities \citep{harris2021moral,bengio2025illusions}. 
A major area of research is to improve the safety on the technical side.
Ensuring the safety of agentic LLMs requires moving beyond prompt-based defenses toward integrated, multi-layer safeguards. In the field of safe reinforcement learning, a key direction is to include explicit safety constraints on the agent's actions and  within the agent's reasoning and planning pipeline \citep{gu2024review}. 

Another approach to improve safety is to employ continual adversarial training and automated red-teaming to enhance robustness against manipulation. Further priorities are developing mechanisms for self-regulation and risk awareness, enabling agents to detect and avoid unsafe actions, and establishing rigorous, standardized safety benchmarks \citep{andriushchenko2024agentharm}. Together, these measures outline a roadmap for accountable and trustworthy deployment of agentic LLMs in high-stakes domains. 

However, legal and ethical questions about who is responsible for the actions of an agent remain. Now that agents are starting to be employed in the real world, addressing these issues becomes of paramount importance. 
Clearly, many safety, ethics and trust issues
will have to be addressed before the full breadth of the possibilities
of agentic LLMs can be enjoyed.
Safety and ethics will become important topics on the research agenda of agentic LLMs, deserving their own surveys and books \citep{gan2024navigating,jiao2025navigating,raza2025responsible}.




\subsection{Crime}
The discussion on ethics and consciousness becomes quite concrete when we consider criminal behavior. Technological developments can be used by people who act within the law, and those who act outside of the law. They can also be used for unethical behavior for which no effective laws yet exist. 
As agentic LLMs are able to display more and more believable behavior, it will become difficult to distinguish real and artificial behavior. This may lead to dangerous ethical situations, for example when criminals use the technology to influence the behavior of unsuspecting victims.

The line between lawful persuasion, misuse of confidence, and fraud, can be unclear. 
Agentic LLMs will greatly increase the sophistication of social engineering using  persuasive algorithms in social networks, marketing, espionage, and cybercrime \citep{kshetri2025transforming,lupinacci2025dark}. Although agentic LLM can also be used to combat crime \citep{challita2025redteamllm,andriushchenko2024agentharm}, it will be some time before societal awareness will have caught up with the risks. %\todo{Lastig om goede statistieken te vinden over cybercrime}
Especially people who are less familiar with social networks and LLMs may be susceptible to confidence tricks and agentic LLM cybercrime. 


\subsection{Limitations of Applications}
The list of applications of agentic LLMs is large. Proofs of concept have been created that impress us with their results, at first sight. When we look deeper at the results, limitations appear that may not be easy to fix and that may have serious consequences when the agents are used in the real world. 

\subsubsection{AI Science}
Many agentic LLMs are developed to support scientific research. All phases of the research pipeline are being supported. Tools for summarizing literature helps the ideation phase, tools exist to generate scripts for running experiments in different fields, there is support for data interpretation and for creation of graphs, writing support exists, as well as tools to support the reviewing of scientific manuscripts.  Indeed, it is becoming easier and easier to produce research papers that are well-written and that look good, on the surface. 

In mathematics, impressive results are reported that LLM agents perform on the level of math olympiads, both in geometry \citep{trinh2024solving} and in algebra  \citep{hubert2025olympiad}. These systems include transformer networks, but are otherwise  designed as mathematics solvers, using specialized automated theorem proving technology such as Lean \citep{de2015lean} and reinforcement learning \citep{silver2017mastering,browne2012survey}.

In 2025 a benchmark study was performed on  AI science tools on a range of domains from machine learning \citep{nathani2025mlgym}. Using  then-current frontier models, the authors  did not find that models could generate novel
hypotheses, algorithms, architectures, or make substantial scientific improvements. In particular, the authors note that {\em modern LLM agents can successfully tackle a diverse array of quantitative
experiments, reflecting advanced skills and domain adaptability} but also that {\em it is not yet clear if the notion of scientific novelty can
be successfully automated or even formally defined in a form suitable for agents}.



\subsubsection{E-commerce}
Liability for actions is a major problem for deploying agentic LLMs in e-commerce. Security and privacy concerns will inhibit users from  providing personal details and preference to leaky computer systems. This leads to the catch-22 situation of recommender systems: when users are unwilling to provide elaborate information about their true preferences, the quality of suggestions will be low, users will have low trust in the system, and will be even less willing to provide information. Only when it is easy to collect high quality information about preferences, in a safe and secure way, and when the models are able to provide strong and useful suggestions, will agentic LLM for e-commerce be adopted by the public.

\subsubsection{Trading}
For financial trading liability is of paramount  importance. 
Studies report that LLMs can be better than human analysts in summarizing large volumes of financial reports. However, in financial trading the stakes are high. Wrong trading advice will directly lead to financial loss. Furthermore, given the unpredictable nature of financial markets, it may take some time to recognize an erroneous model, since it will not be clear when an individual loss is caused be an imperfection in the model or just by an unlucky trade. Models should be thoroughly tested before they will be trusted enough to be used. Unless bulletproof, automated trading systems will not go beyond support of human decision making.

LLMs in finance are mostly successful in analyzing textual reports. For time series analysis of financial data some results are published, but it is hard to achieve reliable and verifiable benchmark results.
Although some studies claim good results, others report that predictions by agentic LLM of financial market results remain in doubt \citep{li2026can}. %are better thanSummary of text, yes, quality of quantitative predictions low. 

 


\subsubsection{Medicine}
Accuracy and especially explainability of results are even more important for trustworthiness in medicine. In medicine, model predictions are directly related  to liability and ethics of actions. For medical applications, explainability of results is essential, and agentic LLM systems will include a human in the decision making loop.

As in trading, the most promising results of agentic LLM in medicine are in the analysis and summarization of textual documents, such as meeting reports or literature surveys. Results in other applications, such as decisions on treatment plans, or performing operations, are limited. 

Medicine might well be one of the most valuable and at the same time one of the most conservative fields of application for agentic LLM. In matters of life and death, we should be extra  careful, and at the same time the potential gains are large.



\section{Concluding Remarks}
%Some appropriate words on the future of agentic LLM. 

Agentic LLM technology is developing at a rapid pace. This book has 
traced that development across three interconnected pillars: \emph{reasoning}, 
\emph{acting}, and \emph{interacting}. Reasoning LLMs provide the cognitive 
foundation, through chain of thought, self-reflection, and retrieval 
augmentation: models have become capable of multi-step problem solving, 
literature synthesis, and complex decision support. Acting agents extend this 
foundation into the real world, using tools, APIs, and world models to assist 
users in domains from robotics to finance to medicine. Interacting agents go 
further still, engaging with one another in multi-agent simulations that 
exhibit emergent social behaviors such as cooperation, trust, and norm formation, that 
were not explicitly scripted in the prompts.

These three pillars are connected. As Figure~\ref{fig:virtuous} 
illustrates, they form a  cycle: reasoning enables better acting and 
interacting, while the data generated by (inter)acting  agents can be used to pretrain and finetune better reasoning models. This virtuous cycle 
offers a promising path beyond plateaus in LLM scaling, as 
inference-time interaction data from agentic assistants, robotic systems, and 
multi-agent simulations, provides a continuously replenishing source of 
high-quality, diverse training signals.

Yet, significant challenges remain across all three pillars. Hallucination 
undermines trust in reasoning chains and threatens the stability of 
self-improving systems. Safety, liability, and ethical accountability are 
unresolved for acting agents in high-stakes domains. Emergent behavior in 
large-scale multi-agent systems remains difficult to validate and control. 
Further research in the fields of mechanistic interpretability, open world models, scalable simulation infrastructures, and robust safety 
benchmarks is required to overcome these challenges.
 
We are at an early and consequential moment. Proofs of concept are 
impressive; real-world deployments are beginning; and the theoretical 
foundations---of memory, metacognition, consciousness, and collective 
intelligence---remain open questions that connect artificial intelligence to 
philosophy, cognitive science and law. This book has aimed to give you the 
technical grounding and conceptual vocabulary to contribute to these 
questions. The field needs researchers who understand both the power and the 
limits of what agentic LLMs can do today, and who can help shape what they 
will become tomorrow.


\begin{comment}

\chapter{Conclusion/Future}\label{chap:conc}
The interest of society and the AI research community  in Agentic LLM is
large, and many papers have been written.  We have reviewed the field, with an emphasis on the most recent works.

%\paragraph{For Society}
There is an interest from society in Agentic LLM. LLM agents can  assist us in our daily
lives in many ways---from writing essays, booking flights, having pleasant and
interesting conversations, folding our laundry,
to making better medical diagnoses, performing
better stock analyses, to
support healthy lifestyle changes, to make sure we take our medicine, to assist us when we are less mobile. Tool use by assistants is enabled by technology from the first category of the taxonomy: reasoning LLMs, self reflection and retrieval augmentation. 

%\paragraph{For Science}
There is also an %great 
interest from science in Agentic LLM. %multi-agent simulation. %from the social sciences, including economics and political science
% in multi agent LLM simulations.
%: when LLMs become agents; with
Since LLM agents can now  interact in natural language, agent behavior can be better understood, and multi-agent simulations can be made more realistic more easily than before.
%whose
%behavior becomes more realistic, 
Important questions in social and political science can be
researched, such as in game
theory (social dilemmas), social interaction (negotiation, theory of
mind),  and societal development
(cooperation, norms,  the rise and fall of empires).
Some of these  goals are  within  reach,  some have been realized already, and some become a possibility. 
%From our survey, we conclude that the commercial and scientific expectations
%are warranted. We also believe that many more fields (chemistry?) will
%find uses for the progress afforded by the combination of agent
%technology with LLM technology.
Agent interactions are enabled by the previous two categories (reasoning and acting): social reasoning benefits from self reflection, negotiation actions are grounded,  information can be retrieved.

%\paragraph{For Artificial Intelligence}
Assistant interactions and agent simulations generate new training data. Where %There is also great interest from the LLM community  in Agentic LLMs. 
there are signs that LLM training on existing data from the
internet is plateauing, the newly generated interaction data can be used for further training. This data is not human-generated, but from interactions with artificial agents. It remains  a topic of future research to understand how this impacts LLM training and finetuning. Results with Vision-Language-Action models that are trained on self-generated action sequences % of actions that are generated by robotic interaction 
are promising. They show results in domestic tasks (kitchen tasks, folding laundry) that had not been achieved by other machine learning methods.

%. Agentic LLM can be a  way to generate new, high quality, interactions, and thus new training data. We note, however, that many methods that are inference time methods, and the training scheme used is in-context learning, where model parameters are not changed. In order to train the LLM with the new data that is generated by the agentic methods, the data should be used in finetuning and pretraining schemes.

%\subsection{Advantages}
%The joining of LLMs and agents has opened up many opportunities and goals.
%To reach these goals, the papers that we have reviewed aim for LLM agents
%that are (1) intelligent, (2) active, and (3) cooperative, for the
%good of society and to generate more training data for LLMs.

%We have reviewed papers that report significant progress in all three fields.
The use of reasoning models for areas such as shopping recommendations, financial analysis, and medical diagnosis has attracted attention to AI assistants. LLM agents can also navigate social dilemmas and realistically simulate emergent behavior in multi-agent interactions. % prompted research in LLM-based agent simulations. 
Reinforcement learning results in robotics with Vision-Language-Action models have attracted interest in inference time data to augment finetuning. Together, these developments have attracted interest in Agentic LLMs, and the synergy between  three developments has increased the interest: reasoning LLMs allow strong decision making; agent technology allows assistants
to interact in the real world; these interacting agents generate behavior on which models can be pretrained and finetuned, improving the performance of the base LLMs. Figure~\ref{fig:virtuous} illustrates this cycle of continuous improvement. 


%\subsection{Challenges}
%Agentic LLMs are at the beginning of their development, with many
%experiments of new ideas still being tried. 
%The field is quite active, the interest of research is high, 
%as are commercial expectations. Given how young the field is, there is a tension between the level of expectations of society and the level of maturity of research results. 
%Many challenges still remain.

%\subsubsection{Reasoning}
%We will list challenges to improve step by step reasoning intelligence of Agentic LLMs
%further.
\subsection{Research Agenda for Agentic LLM}
The  interest from researchers in Agentic LLMs is large. Our survey has yielded interesting topics for a research agenda for Agentic LLM. Please refer to Table~\ref{tab:agenda} for a summary of the agenda.


\paragraph{Hallucination and Stable  Behavior}
A challenge for the virtuous autocurriculum cycle is that LLMs hallucinate. LLM answers may look good, but be factually wrong. Reasoning chains may be unfaithful, giving good answers for the wrong reason, and giving wrong answers when least expected. Especially when such dubious results are used to in further train the LLM, this training may diverge. In social simulations, emerging behavior patterns, such as cooperation, fairness, trust or norms, may collapse. In multi-step reasoning, self verification and self consistency methods were developed to address error accumulation. In reinforcement learning, exploration and diversity are important methods to ensure good coverage of the state space. 
In social simulations and gaming, open world models and open-ended behavior are being used to stimulate exploration and diversity. Such models can provide suitable environments for automated generation of training curricula. 

Faithfulness for Chain of Thought is studied by
\citet{lyu2023faithful,lanham2023measuring,turpin2024language}.
Mechanistic interpretability can provide ways to look inside the LLM,
to better understand if the model follows the reasoning steps that we expect it to take
\citep{nanda2023progress,bereska2024mechanistic,ferrando2024primer}. 
The conditions that influence stability of emergent behavior (cooperation, fairness, trust) may be studied further.

For Agentic LLMs that learn  from their own results, other methods will need to be developed, and hallucination features prominently on the research agenda for Agentic LLM, with mechanistic interpretability and open world models and norms as important items.

% In step-by-step reasoning it leads to faithfulness issues, in assistants it leads to safety and security issues, in multi agent simulations it leads to low robustness in understanding social dilemmas. Research into Mechanistic Interpretability \citep{nanda2023progress} may yield more insights into what is going on inside the models, and may yield ways to reduce hallucination.  Open world environments may help provide diverse training data for hardening models, reducing hallucination.

\paragraph{Agent Behavior at Scale}
Studies of emergent behavior need realistic agent behavior, and we expect more research to be performed to improve agent behavior, for example by closely modeling human behavior in generative agents \citep{park2024generative}. Some behavior patterns in multi agent simulations emerge at scale, as  studies with specialized agent infrastructures have shown \citep{park2023generative,yang2024oasis,wu2023autogen}. However, the  number of LLM agents that can be simulated reliably is often limited. Although role-based agent approaches may improve scalability, we believe that  more research into scaling of simulations with complex agents is necessary.

Related to this challenge is the cost of training LLMs. Pretraining and finetuning an LLM can be expensive. Knowledge distillation is  a popular method to extract essential knowledge and behavior from a large model into a small model, at much lower computational cost \citep{xu2024survey}.  Experiments have
shown that reasoning steps can be distilled from large to smaller
language models \citep{gu2023minillm,li2023symbolic,muennighoff2025s1}. Knowledge distillation in LLM agents in an important item for our research agenda. 

The debate on artificial super-intelligence is fueled, in part, by the growing performance of individual LLMs, which is an important aspect of Agentic LLM research. Another aspect of Agentic LLM research is the study  of emergent behavior, of cooperation and trust in agentic societies.  Such studies of agent behavior at scale  may influence our view on the nature of intelligence, and may influence the
discussion on artificial super-intelligence and the future of  society. Furthermore, the world around us is organized in groups in which power hierarchies
are prevalent. Many multi agent simulations assume a flat power
hierarchy. %More research is needed into power dynamics.
Multi agent simulations should also go beyond equality.


\paragraph{Self Reflection}
%\paragraph{Self Reflection}
Self reflection mechanisms are used in  advanced prompt-improvement algorithms. External explicit prompt management may be error prone and brittle. An alternative is  to let the LLM perform the self reflection and step-by-step management internally, as in the original  Chain of Thought. 
DeepSeek R1 \citep{guo2025deepseek} is a reasoning model that is trained (finetuned) by the GRPO reinforcement learning method \citep{shao2024deepseekmath}. The model is trained  on its own reasoning results, and was found to self-reflectively reason over its own results, identifying effective reasoning patterns implicitly. 

\citet{schultz2024mastering} train a model on search sequences \citep{gandhi2024stream} in games such as chess, and VLAs are trained on action sequences \citep{kim2024openvla}. Their work shows that, in addition to implicit step-by-step reasoning, implicit search is possible. An open question is whether LLMs can perform self reflection internally. %Can we find ways to
%follow tree shaped search paths  implicitly, with a single prompt,
%without the need for an external algorithm?  

The addition of external state to an LLM has enabled reasoning and a form of self
reflection, which is a form of metacognition (thinking about thinking). 
%Some of the prompt-improvers exhibit a form of self reflection. Can
%LLM agents also exhibit meta cognition? How can we test if this has been achieved?
%Advanced prompt-improvement algorithms are typically external
%optimization or tree traversal loops. Simple step-by-step sequences are
%implicit (single prompt). For the sake of elegance, can we get the LLM to do advanced
%optimization/tree traversal with a single prompt, implicitly? And can we then use 
%mechanistic interpretability to verify that the model  follows a tree?
%
%
%\paragraph{Meta Cognition}
% 
LLMs that reflect on their own behavior raises visions of true
artificial intelligence. %Meta cognition is thinking about
%thinking. 
If  LLMs can self-reflect, can they exhibit metacognition \citep{wang2023metacognitive,didolkar2024metacognitive}? 
Self reflection by LLMs is another item for the research agenda. 
% How can we test this? 

Self reflective methods are being used to create agents to perform scientific discovery \citep{eger2025transforming}. How these agents will influence, and possibly improve, the process of scientific discovery is an exciting area of research. 

%\paragraph{Identity}
When we add outside state to the input prompts, the input to the LLM will differ based on the history, and so will the answers of the LLM. Differences in memory may be preceived as a
 personality of the LLMs by its users. The question if LLMs with outside memory exhibit a personality is a topic for future research. 
%Does the personality influence
%the behavior of the LLM? Do humans recognize different LLM
%personalities? Do LLMs have an identity? Should agents have rights?



%\subsubsection{Active}
%We will list the challenges to improve Agentic LLM actions in the
%world further.


\paragraph{Safety}
Safety is an important issue in LLMs that act in the world. The
problem is  studied, but far from being solved
\citep{brunke2022safe,andriushchenko2024agentharm,samvelyan2024rainbow}.
Actions by assistants and robots in the real world have real world consequences. When a financial trading assistant hallucinates, or when a self driving robot makes a wrong inference, questions on responsibility and liability should be addressed. More legal and ethical questions arise, for example, on privacy and fairness, and, possibly, concerning the rights of algorithmic entities \citep{harris2021moral}. 

The application areas  for the assistants in this survey---shopping, medical diagnosis, finance---are  narrow. The narrower the application domain, the better the answers. 
Clearly, many safety, ethics and trust issues
will have to be addressed before the full breadth of the possibilities
of Agentic LLMs can be enjoyed.
Safety will continue to be an important topic on the research agenda of Agentic LLM.


\begin{table}
    \centering%\footnotesize
    \begin{tabular}{ll}
     {\em Topic} & {\em Challenge} \\ \hline 
        Training Data & Finetune with inference time reasoning data\\
        & Convergent/stable reinforcement learning  \\
        & VLA, Multimodal signals, such as speech \\ \hline
         Hallucination & Use Self Verification \\
         & Use Mechanistic Interpretability \\
         & Use Open Ended/Open World Models for exploration\\ \hline
         Agent Behavior & Scalable simulation infrastructure, role playing\\
         & Distill reasoning to small models\\
         & Models of agent and human behavior, emergent behavior,  future of society\\ \hline
         Self reflection & In-model self reflection and metareasoning\\
         & Metacognition, personality\\ 
         & Automated Scientific Discovery \\ \hline
         Safety & Assistants: Responsibility, liability\\
         & Privacy, fairness of data\\
         & Wider  application areas for assistants\\ \hline
    \end{tabular}
    \caption{Summary of Research Agenda for Agentic LLM}
    \label{tab:agenda}
\end{table}

%\begin{table}
%    \centering\footnotesize
%    \begin{tabular}{lll}
%    {\em Area} & {\em Topic} & {\em Challenge} \\ \hline 
%        Intelligent & Multi Step Reasoning %Inference-time data   
%        & Use inference-time results to finetune models  \\
%        &Self Reflection & Explore Implicit search/Meta-cognition  \\
%        & Retrieval Augmentation & Distillation to Small Language Models \\ \hline
%    
%        Active &  Robot/Tools %Hallucination
%        & Hallucination: Mechanistic Interpretability   \\ 
%        & & Develop Reflection/Verification methods \\ 
%        & & Use Open World agents/Automated curricula \\ 
%         & Action Models & Improve VLAs/World Models \\
%         & Assistants & Study privacy and  liability issues \\ \hline
%        LLM-Agent Based Models & Dyadic Interactions %Robust agent behavior  
%        & Improve modeling with agent roles \\ 
%         & Agent Societies % Robust emergent behavior  
%                                & Develop scaling infrastructures\\ 
%        & Small Groups & Explore power hierarchy models \\  \hline
%    \end{tabular}
%    \caption{Summary of Research Agenda Agentic LLM}
%    \label{tab:agenda}
%\end{table}

\section{conclusion}
We survey recent papers on Agentic LLMs, with a taxonomy based on three categories: reasoning/acting/interacting. The first part of the survey covers approaches on {\em reasoning} and retrieval.
Initially, LLMs score high on language understanding tasks such as general question answering, but not on reasoning tasks such as math word problems. Using approaches based on Chain of Thought's {\em Let's think step by step} prompt, reasoning LLMs have become better in solving math word problems, and also in solving more complex reasoning tasks. Models are also surprisingly good at analyzing medical diagnosis documents, financial statements, and social dilemmas. The first part of the survey covers the reasoning technology that enables the next two categories, where agents  act and interact.

LLMs have been enhanced with interfaces to use tools for information retrieval, and to act in the real world, potentially making capable assistants. Reasoning and self reflection allow for  complex interactions and self verification. The second part of the survey covers assistants that {\em act} (including robots and world models). Enhancing LLMs with action tools (such as for APIs) turns them into agents. Examples of such agents are assistants that analyze medical diagnosis documents, or that recommend trading actions on the stock market. Although language can have real world consequences,  the surveyed papers describe proofs of concept, but not yet actual deployments in the real world. Results are promising, but before real world deployment can be considered, questions on liability may need to be addressed first, especially  in high impact applications in medicine and finance. Applications in companionship, and for shopping recommendations, may be closer to realization.

In addition to the situation where a single assistant-LLM agent interacts with the world, agents can also interact with other agents. The third part of the survey covers {\em interaction}:  multi-agent simulations of societies of agents. LLMs are reasonably good at reasoning about agent behavior (theory of mind), and at understanding social dilemmas such as the prisoners dilemma and public goods games.  In comparison to classical agent-based simulations of social behavior, LLMs have increased the realism of such simulations, both qualitatively (one paper describes a realistic Valentine's day neighborhood party) and quantitatively (one paper reports a simulation of a social network with millions of agents). Social LLM simulations exhibit complex emerging behaviors such as cooperation, trust, and norms. Self reflection can improve LLM reasoning about behavior further.

\paragraph{New data}
At another level, taken together, the three areas of reasoning/acting/interacting enable an important change in the machine learning paradigm for LLMs: Agentic LLMs move from (self-)supervised learning to reinforcement learning.
%
Models learn from datasets---supervised---but agents learn from their own actions---interactively. LLM learning has been reported to plateau, since it becomes harder and harder to find new/better training datasets. The reinforcement learning approach lets agents act in an environment, improving the policy from environment feedback, changing parameters from inference time results. In Agentic LLM the models are active and can learn from self-initiated feedback in assistant and multi-agent simulation settings. They can continue to learn as long as the environments are challenging, just like AlphaGo Zero taught itself to play Go,  by playing against a copy of itself, with zero help from datasets of grandmaster games (and surpassing all human grandmasters). In agentic LLMs, inference time results can be used for pretraining and finetuning. In robotics, promising early results with Vision-Language-Action models are described, and some reasoning LLMs are being trained with reinforcement learning.

\paragraph{Contributions}
The contributions of this survey are as follows: we define Agentic LLMs as reasoning LLMs that act in the world and that interact with other agentic LLMs. In our survey: (1) {\em reasoning} LLMs provide the enabling technology for agentic LLMs; (2) interfaces to tools and APIs allow LLMs to {\em act}, turning them into agents, they can interact with the real world, as assistants for their users; (3) agent LLMs can also {\em interact} with each other, reflect on interaction  of self and others, in  multi-agent simulations of complex, emergent, cooperative, social behavior. 

Our three-way taxonomy for agency is built on the reasoning technology of category 1,  which is growing in importance in agentic LLMs. Robotic interaction and tool-use builds on grounded retrieval augmentation, social interaction (such as theory of mind) builds on self reflection, and all categories benefit from reasoning and self-verification. Closing the cycle, the acting and interacting categories generate training data for further pretraining and finetuning LLMs,  beyond plateauing traditional datasets (Figure~\ref{fig:virtuous}). The impressive generalization capabilities of Vision-Language-Action models are testament to the power of this approach. 


The reasoning paradigm connects to works in human cognition, and some papers anthropomorphize LLM computations in Kahneman's terms of System 1 thinking (fast, associative) and System 2 thinking (slow, deliberative). Works on reasoning focus  on intelligence of single (individual) LLMs. This individualistic view also gives rise to  discussions about superintelligence, some utopian, and some not. 

The agentic paradigm enables an element of machine learning that is new for LLMs. In reinforcement learning, agents self reflect and choose their own actions, and learn from the feedback of the world in which they operate.   No dataset is needed, nor is learning limited by the complexity of a dataset, it is only limited by the complexity of the world around the agent. The agent paradigm  creates a more challenging training setting, allowing Agentic LLMs to keep on improving themselves, not being held back by human datasets.

The multi-agent paradigm studies agent-agent societies. The focus is on emergent behaviors such as egoism/altruism, competition/collaboration, and (dis)trust.   Social cognitive  development and the emergence of collective intelligence are also studied in this field. Connecting back to the reasoning paradigm, multi-agent studies may inform discussions about (super)intelligence. 

Our  agenda for future research includes suggestions for all three aspects of the taxonomy. In reasoning, many challenges remain in self reflection, metacognition and implicit search methods. The stability of training for self-generating inference time data should improve, for example with methods to improve exploration with open world models, and with mechanistic interpretability methods to better understand hallucination and faithful reasoning. For assistants, further work is required to improve performance in a wider application portfolio, as well as on ethical and legal issues. Agent simulations will yield better models of human and group behavior; we suggests work on improving  modeling infrastructures and on improving the realism of agent behavior. Furthermore, energy consumption may be reduced with small language models, using distillation methods; multi-agent studies may teach us about social behavior of LLM-agents (AI4Good). 

\end{comment}

\bibliographystyle{plainnat}
\bibliography{agenticbook,refs,case,synthetic}
